% ============================================================
%  Chapitre 1 — Introduction et Statistiques Descriptives Multivariées
% ============================================================
\section{Introduction et Statistiques Descriptives Multivariées}

% ────────────────────────────────────────────────────────────
\subsection{1.1 Position de l'analyse des données}
% ────────────────────────────────────────────────────────────
\begin{frame}{1.1 — Qu'est-ce que l'analyse des données ?}
\begin{columns}[T]
\column{0.55\textwidth}
\begin{definition}
L'\textbf{analyse des données} (\textit{data analysis}) est l'ensemble des
techniques exploratoires \alert{descriptives}, fondées sur la
\textbf{décomposition matricielle} et la \textbf{géométrie des nuages de points},
visant à révéler la structure d'un tableau de données
$\mathbf{X} \in \R^{n \times p}$ \emph{sans modèle probabiliste}.
\end{definition}

\vspace{0.3cm}
\textbf{Tradition française} (Benzécri, années 1960) :
\begin{itemize}\small
  \item \alert{Géométrique} (inertie, projections orthogonales)
  \item \alert{Descriptive} (pas d'hypothèses inférentielles)
  \item \alert{Visuelle} (cartes factorielles)
\end{itemize}
\column{0.42\textwidth}
\begin{tikzpicture}[scale=0.85, font=\scriptsize,
  box/.style={draw=deepblue, fill=deepblue!10, rounded corners,
              minimum width=3.3cm, minimum height=0.65cm, align=center}]
  \node[box]                      (stat) at (0, 2.2) {Statistique\\inférentielle};
  \node[box, fill=deeporange!10, draw=deeporange]  (ad)   at (0, 0.9) {Analyse\\des données};
  \node[box, fill=darkgreen!10, draw=darkgreen]    (ml)   at (0,-0.4) {Machine\\Learning};
  \node[font=\tiny, gray, align=center] at (4.4, 2.2) {modèle\\$+$ tests};
  \node[font=\tiny, gray, align=center] at (4.4, 0.9) {description\\géométrique};
  \node[font=\tiny, gray, align=center] at (4.4,-0.4) {prédiction\\$+$ optimisation};
  \draw[->, thick, gray] (stat.south) -- (ad.north);
  \draw[->, thick, gray] (ad.south)   -- (ml.north);
\end{tikzpicture}
\end{columns}
\end{frame}

% ────────────────────────────────────────────────────────────
\begin{frame}{1.1 — Objectifs de l'analyse des données}
\begin{columns}[T]
\column{0.49\textwidth}
\textbf{Trois grands objectifs}
\begin{enumerate}
  \item \textbf{Réduire la dimension} : remplacer $p$ variables corrélées
        par $q \ll p$ \emph{composantes synthétiques}.
  \item \textbf{Représenter graphiquement} les individus et les variables
        sur des plans factoriels.
  \item \textbf{Classifier} (typer) les individus en groupes homogènes.
\end{enumerate}

\vspace{0.3cm}
\textbf{Méthodes selon le type de données}
\begin{itemize}\small
  \item Quantitatives $\to$ \alert{ACP}
  \item Tableau de contingence $\to$ \alert{AFC}
  \item Qualitatives $\to$ \alert{ACM}
  \item Mixtes $\to$ AFDM, AFM
\end{itemize}
\column{0.49\textwidth}
\begin{block}{Place de ce cours}
  Articulation entre :
  \begin{itemize}\small
    \item \textbf{algèbre linéaire} (LMAD, théorèmes)
    \item \textbf{statistique multivariée} (GAMMA)
    \item \textbf{interprétation métier} (BA)
  \end{itemize}
\end{block}

\begin{alertblock}{Différence avec le ML supervisé}
  Pas de variable cible $y$ : on cherche à \alert{décrire} et
  \alert{comprendre} la structure, pas à \alert{prédire}.
\end{alertblock}
\end{columns}
\end{frame}

% ────────────────────────────────────────────────────────────
\subsection{1.2 Tableau de données}
% ────────────────────────────────────────────────────────────
\begin{frame}{1.2 — Le tableau de données}
\begin{definition}[Tableau de données]
On note $\mathbf{X} \in \R^{n \times p}$ le tableau dont la ligne $i$ (individu) est
$\mathbf{x}_i = (x_{i1}, \ldots, x_{ip})^\top \in \R^p$, et la colonne $j$
(variable) est $\mathbf{x}^j = (x_{1j}, \ldots, x_{nj})^\top \in \R^n$.
\end{definition}

\begin{columns}[T]
\column{0.50\textwidth}
\textbf{Deux lectures}
\begin{itemize}
  \item \textbf{Ligne} : nuage des $n$ individus dans $\R^p$
  \item \textbf{Colonne} : nuage des $p$ variables dans $\R^n$
\end{itemize}
\vspace{0.2cm}
\textbf{Types de variables}
\begin{itemize}\small
  \item \emph{Quantitatives} (continues, discrètes) : ACP
  \item \emph{Qualitatives nominales} : ACM
  \item \emph{Qualitatives ordinales} : recodage prudent
  \item \emph{Effectifs / contingence} : AFC
\end{itemize}
\column{0.46\textwidth}
\centering
\footnotesize
\renewcommand{\arraystretch}{1.1}
\begin{tabular}{c|ccccc}
  & $V_1$ & $V_2$ & $\cdots$ & $V_p$ \\
\hline
$i_1$    & $x_{11}$ & $x_{12}$ & $\cdots$ & $x_{1p}$ \\
$i_2$    & $x_{21}$ & $x_{22}$ & $\cdots$ & $x_{2p}$ \\
$\vdots$ & $\vdots$ & $\vdots$ & $\ddots$ & $\vdots$ \\
$i_n$    & $x_{n1}$ & $x_{n2}$ & $\cdots$ & $x_{np}$ \\
\end{tabular}

\vspace{0.3cm}
\begin{block}{\small Exemple}
\small Tableau \texttt{decathlon} : $n=41$ athlètes,
$p=10$ disciplines $+$ classement.
\end{block}
\end{columns}
\end{frame}

% ────────────────────────────────────────────────────────────
\subsection{1.3 Statistiques descriptives univariées}
% ────────────────────────────────────────────────────────────
\begin{frame}{1.3 — Caractéristiques d'une variable}
\begin{columns}[T]
\column{0.50\textwidth}
\textbf{Tendance centrale}
\begin{align*}
\bar{x}^j &= \frac{1}{n}\sum_{i=1}^n x_{ij}
  \quad\text{(moyenne arithmétique)}\\
\tilde{x}^j &= \text{médiane}(\mathbf{x}^j)
\end{align*}

\textbf{Dispersion}
\begin{align*}
s_j^2 &= \frac{1}{n}\sum_{i=1}^n (x_{ij} - \bar{x}^j)^2
  \quad\text{(variance empirique)}\\
s_j   &= \sqrt{s_j^2} \quad\text{(écart-type)}
\end{align*}

\begin{rema}
On utilise ici la variance \alert{biaisée} (divisée par $n$),
adaptée à l'analyse géométrique.
\end{rema}
\column{0.46\textwidth}
\textbf{Quantiles et étendue}
\begin{itemize}\small
  \item $Q_1, Q_2 = \tilde{x}, Q_3$ : quartiles
  \item $IQR = Q_3 - Q_1$ : écart interquartile
  \item Étendue : $\max - \min$
\end{itemize}

\textbf{Forme}
\begin{itemize}\small
  \item Asymétrie (skewness) : $\gamma_1 = \E[(X-\mu)^3]/\sigma^3$
  \item Aplatissement (kurtosis) : $\gamma_2 = \E[(X-\mu)^4]/\sigma^4 - 3$
\end{itemize}

\begin{alertblock}{Présence d'outliers}
  Préférer $\tilde{x}$ et $IQR$ à $\bar{x}$ et $s$ pour résumer une
  variable contaminée.
\end{alertblock}
\end{columns}
\end{frame}

% ────────────────────────────────────────────────────────────
\subsection{1.4 Statistiques descriptives bivariées}
% ────────────────────────────────────────────────────────────
\begin{frame}{1.4 — Liaison entre deux variables}
\begin{definition}[Covariance et corrélation empiriques]
Pour deux variables $\mathbf{x}^j, \mathbf{x}^k \in \R^n$ :
\[
\Cov(\mathbf{x}^j, \mathbf{x}^k)
= \frac{1}{n}\sum_{i=1}^n (x_{ij}-\bar{x}^j)(x_{ik}-\bar{x}^k),
\qquad
\Cor(\mathbf{x}^j, \mathbf{x}^k) = \frac{\Cov(\mathbf{x}^j, \mathbf{x}^k)}{s_j \, s_k}.
\]
\end{definition}

\begin{prop}[Propriétés de la corrélation]
\begin{enumerate}
  \item $\Cor(\mathbf{x}^j, \mathbf{x}^k) \in [-1, 1]$ \hfill (Cauchy-Schwarz)
  \item $|\Cor| = 1 \iff \mathbf{x}^k$ est fonction \emph{affine} de $\mathbf{x}^j$
  \item $\Cor$ est invariante par changement d'unités linéaire (échelle, origine)
  \item $\Cor = 0 \not\Rightarrow$ indépendance (mesure uniquement la liaison \emph{linéaire})
\end{enumerate}
\end{prop}

\vspace{0.1cm}
\begin{block}{Forme matricielle}
$\mathbf{V} = \frac{1}{n} \mathbf{X}_c^\top \mathbf{X}_c$ : matrice de variance-covariance ($p\times p$).\\
$\mathbf{R} = \mathbf{D}_{1/s}\, \mathbf{V}\, \mathbf{D}_{1/s}$ : matrice de corrélation, où $\mathbf{D}_{1/s} = \text{diag}(1/s_j)$.
\end{block}
\end{frame}

% ────────────────────────────────────────────────────────────
\subsection{1.5 Centrage et réduction}
% ────────────────────────────────────────────────────────────
\begin{frame}{1.5 — Centrage et réduction}
\begin{columns}[T]
\column{0.52\textwidth}
\textbf{Centrage}
\[
\mathbf{X}_c = \mathbf{X} - \1_n \bar{\mathbf{x}}^\top, \quad
\bar{\mathbf{x}} = \frac{1}{n}\mathbf{X}^\top \1_n
\]
$\Leftrightarrow$ déplacer l'origine au centre de gravité $\mathbf{g}$.

\vspace{0.3cm}
\textbf{Réduction (standardisation)}
\[
\mathbf{Z} = \mathbf{X}_c \, \mathbf{D}_{1/s}, \quad z_{ij} = \frac{x_{ij} - \bar{x}^j}{s_j}
\]
$\Leftrightarrow$ rendre les variables \alert{adimensionnelles}.

\begin{theorem}[Effet sur la variance-covariance]
\[
\mathbf{V}(\mathbf{X}_c) = \mathbf{V}(\mathbf{X}), \qquad
\mathbf{V}(\mathbf{Z}) = \mathbf{R}(\mathbf{X}).
\]
\end{theorem}
\column{0.44\textwidth}
\begin{block}{Pourquoi réduire ?}
\begin{itemize}\small
  \item Variables d'\alert{unités différentes} (kg, EUR, années)
  \item Variances très différentes
  \item Sans réduction, l'ACP est dominée par la variable de plus grande variance
\end{itemize}
\end{block}

\begin{alertblock}{Quand ne pas réduire ?}
\small Lorsque les variables sont
\textbf{de même nature} et de variances comparables
(ex : températures aux mois de l'année).
\end{alertblock}
\end{columns}
\end{frame}

% ────────────────────────────────────────────────────────────
\subsection{1.6 Distance et métrique}
% ────────────────────────────────────────────────────────────
\begin{frame}{1.6 — Distance entre individus}
\begin{definition}[Métrique]
Une \textbf{métrique} sur $\R^p$ est une matrice $\mathbf{M}$ \alert{symétrique
définie positive}. La distance associée entre $\mathbf{x}_i$ et $\mathbf{x}_{i'}$ est :
\[
d_{\mathbf{M}}^2(\mathbf{x}_i, \mathbf{x}_{i'}) = (\mathbf{x}_i - \mathbf{x}_{i'})^\top \mathbf{M} (\mathbf{x}_i - \mathbf{x}_{i'}).
\]
\end{definition}

\begin{columns}[T]
\column{0.49\textwidth}
\textbf{Cas particuliers}
\begin{itemize}\small
  \item $\mathbf{M} = \mathbf{I}_p$ : distance \alert{euclidienne classique}\\
        $d^2 = \sum_j (x_{ij}-x_{i'j})^2$
  \item $\mathbf{M} = \mathbf{D}_{1/s^2}$ : distance \alert{euclidienne réduite}\\
        $d^2 = \sum_j \frac{(x_{ij}-x_{i'j})^2}{s_j^2}$
  \item $\mathbf{M} = \mathbf{V}^{-1}$ : \alert{Mahalanobis}
  \item $\mathbf{M} = \mathbf{D}_{1/x_{\bullet j}}$ : $\chi^2$ (AFC)
\end{itemize}
\column{0.49\textwidth}
\begin{prop}[Choix de la métrique]
Le choix de $\mathbf{M}$ \emph{définit} l'analyse :
\begin{itemize}\small
  \item ACP normée $\equiv$ ACP centrée avec $\mathbf{M} = \mathbf{D}_{1/s^2}$
  \item AFC : tableau de profils, $\mathbf{M}_{\chi^2}$
  \item AFD : $\mathbf{M} = \mathbf{V}_{\text{intra}}^{-1}$
\end{itemize}
\end{prop}

\begin{block}{Pondération des individus}
\small Chaque individu a un poids $p_i \geq 0$ tel que
$\sum_i p_i = 1$. Cas usuel : $p_i = 1/n$.
\end{block}
\end{columns}
\end{frame}

% ────────────────────────────────────────────────────────────
\subsection{1.7 Notion d'inertie}
% ────────────────────────────────────────────────────────────
\begin{frame}{1.7 — Inertie : la mesure de l'analyse des données}
\begin{definition}[Inertie totale d'un nuage]
Soit $\mathcal{N} = \{(\mathbf{x}_i, p_i)\}_{i=1}^n$ un nuage pondéré de centre de gravité
$\mathbf{g} = \sum_i p_i \mathbf{x}_i$. L'\textbf{inertie totale} de $\mathcal{N}$ par rapport à
$\mathbf{g}$, dans la métrique $\mathbf{M}$, est :
\[
I_{\mathcal{N}} = \sum_{i=1}^n p_i \, d_{\mathbf{M}}^2(\mathbf{x}_i, \mathbf{g})
= \sum_{i=1}^n p_i \, (\mathbf{x}_i - \mathbf{g})^\top \mathbf{M} (\mathbf{x}_i - \mathbf{g}).
\]
\end{definition}

\begin{prop}[Lien avec la trace]
Soit $\mathbf{V} = \sum_i p_i (\mathbf{x}_i - \mathbf{g})(\mathbf{x}_i - \mathbf{g})^\top$
la matrice de variance-covariance pondérée. Alors :
\[
\boxed{\,I_{\mathcal{N}} = \Tr(\mathbf{V} \mathbf{M}).\,}
\]
\end{prop}

\begin{rema}
Pour $\mathbf{M} = \mathbf{I}_p$ et $p_i = 1/n$ : $I_{\mathcal{N}} = \Tr(\mathbf{V}) = \sum_j s_j^2$,
soit la \textbf{somme des variances} des variables. En ACP normée :
$I_{\mathcal{N}} = \Tr(\mathbf{R}) = p$.
\end{rema}
\end{frame}

% ────────────────────────────────────────────────────────────
\subsection{1.8 Théorème d'Huygens}
% ────────────────────────────────────────────────────────────
\begin{frame}{1.8 — Théorème de Huygens (décomposition de l'inertie)}
\begin{theorem}[Huygens]
Soit $\mathcal{N}$ un nuage partitionné en $K$ classes $\mathcal{C}_1, \ldots, \mathcal{C}_K$
de poids respectifs $P_k = \sum_{i \in \mathcal{C}_k} p_i$ et de barycentres
$\mathbf{g}_k = \frac{1}{P_k}\sum_{i \in \mathcal{C}_k} p_i \mathbf{x}_i$. Alors :
\[
\boxed{\;
\underbrace{\sum_{i} p_i \, d^2(\mathbf{x}_i, \mathbf{g})}_{I_{\text{totale}}}
=
\underbrace{\sum_{k=1}^K \sum_{i \in \mathcal{C}_k} p_i \, d^2(\mathbf{x}_i, \mathbf{g}_k)}_{I_{\text{intra}}}
\;+\;
\underbrace{\sum_{k=1}^K P_k \, d^2(\mathbf{g}_k, \mathbf{g})}_{I_{\text{inter}}}.
\;}
\]
\end{theorem}

\begin{columns}[T]
\column{0.55\textwidth}
\textbf{Conséquences fondamentales}
\begin{itemize}\small
  \item \alert{Clustering} : minimiser $I_{\text{intra}}$ $\Leftrightarrow$
        maximiser $I_{\text{inter}}$ (k-means)
  \item \alert{Réduction de dimension} : maximiser
        l'inertie projetée $\Leftrightarrow$ minimiser
        l'inertie résiduelle (ACP)
  \item \alert{AFD} : critère de Fisher
        $\max \, I_{\text{inter}} / I_{\text{intra}}$
\end{itemize}
\column{0.42\textwidth}
\begin{tikzpicture}[scale=0.9]
\fill[deepblue!10] (0.5,0.5) circle(0.5cm);
\fill[deeporange!10] (2.5,0.7) circle(0.5cm);
\fill[darkgreen!10] (1.7,2.0) circle(0.45cm);
\foreach \x/\y in {0.2/0.5, 0.4/0.7, 0.7/0.4, 0.5/0.3}
  \fill[deepblue] (\x,\y) circle(1.3pt);
\foreach \x/\y in {2.3/0.6, 2.6/0.8, 2.7/0.5, 2.4/0.9}
  \fill[deeporange] (\x,\y) circle(1.3pt);
\foreach \x/\y in {1.6/2.0, 1.8/2.1, 1.7/1.8}
  \fill[darkgreen] (\x,\y) circle(1.3pt);
\fill[deepblue] (0.45,0.5) circle(2.5pt) node[left, font=\tiny] {$\mathbf{g}_1$};
\fill[deeporange] (2.5,0.7) circle(2.5pt) node[right, font=\tiny] {$\mathbf{g}_2$};
\fill[darkgreen] (1.7,2.0) circle(2.5pt) node[above, font=\tiny] {$\mathbf{g}_3$};
\fill[black] (1.55,1.06) circle(2.5pt) node[right, font=\tiny] {$\mathbf{g}$};
\draw[dashed, gray] (0.45,0.5) -- (1.55,1.06);
\draw[dashed, gray] (2.5,0.7)  -- (1.55,1.06);
\draw[dashed, gray] (1.7,2.0)  -- (1.55,1.06);
\end{tikzpicture}
\end{columns}
\end{frame}

% ────────────────────────────────────────────────────────────
\subsection{1.9 Le tableau comme deux nuages}
% ────────────────────────────────────────────────────────────
\begin{frame}{1.9 — Dualité individus / variables}
\begin{columns}[T]
\column{0.49\textwidth}
\textbf{Nuage des individus $\mathcal{N}_I$ dans $\R^p$}
\begin{itemize}\small
  \item $n$ points $\mathbf{x}_i$
  \item Métrique $\mathbf{M}$ sur $\R^p$
  \item Pondération $\mathbf{D}_p = \text{diag}(p_i)$
  \item Centre $\mathbf{g}$, inertie $\Tr(\mathbf{V}\mathbf{M})$
\end{itemize}

\vspace{0.2cm}
\textbf{Nuage des variables $\mathcal{N}_V$ dans $\R^n$}
\begin{itemize}\small
  \item $p$ points $\mathbf{x}^j$
  \item Métrique $\mathbf{D}_p$ sur $\R^n$
  \item Pondération $\mathbf{M}$ sur les variables
  \item Produit scalaire $\langle \mathbf{x}^j, \mathbf{x}^k \rangle = \Cov(\mathbf{x}^j, \mathbf{x}^k)$
\end{itemize}

\column{0.49\textwidth}
\begin{theorem}[Dualité]
Les deux nuages $\mathcal{N}_I$ et $\mathcal{N}_V$ partagent les
\alert{mêmes valeurs propres non nulles}. La même décomposition
matricielle (SVD de $\mathbf{X}$) fournit simultanément :
\begin{itemize}\small
  \item les axes principaux de $\mathcal{N}_I$
  \item les axes principaux de $\mathcal{N}_V$
\end{itemize}
\end{theorem}

\begin{block}{Conséquence pratique}
\small Les \textbf{individus} et les \textbf{variables} se représentent
sur les \emph{mêmes} plans factoriels (biplot) — l'ACP est
\alert{auto-duale}.
\end{block}
\end{columns}
\end{frame}

% ────────────────────────────────────────────────────────────
\subsection{1.10 Aperçu R}
% ────────────────────────────────────────────────────────────
\begin{frame}[fragile]{1.10 — Aperçu : statistiques descriptives en R}
\begin{lstlisting}
library(FactoMineR)    # methodes factorielles
library(factoextra)    # visualisation
library(corrplot)      # matrice de correlations

data(decathlon)        # 41 athletes x 13 variables
X <- decathlon[, 1:10] # variables actives quantitatives

# Statistiques univariees
summary(X)
sapply(X, sd)

# Centrage / reduction
Z <- scale(X)          # par defaut center=TRUE, scale=TRUE
colMeans(Z); apply(Z, 2, sd)

# Matrice de variance-covariance et de correlation
V <- cov(X)            # divise par n-1 (correction de Bessel)
R <- cor(X)
corrplot(R, method = "color", type = "upper", tl.cex = 0.7)

# Inertie totale (ACP normee)
sum(diag(R))           # = nombre de variables = 10
\end{lstlisting}
\end{frame}

% ────────────────────────────────────────────────────────────
\subsection{1.11 Récapitulatif}
% ────────────────────────────────────────────────────────────
\begin{frame}{1.11 — Récapitulatif du chapitre}
\begin{columns}[T]
\column{0.49\textwidth}
\textbf{Concepts à retenir}
\begin{itemize}\small
  \item Tableau $\mathbf{X} \in \R^{n\times p}$, deux lectures
        (individus / variables)
  \item Matrice de variance-covariance
        $\mathbf{V} = \frac{1}{n}\mathbf{X}_c^\top \mathbf{X}_c$
  \item Matrice de corrélation $\mathbf{R}$
  \item Centrage (décalage), réduction (échelle)
  \item Métrique $\mathbf{M}$ symétrique définie positive
  \item Inertie $I = \Tr(\mathbf{V}\mathbf{M})$
  \item Théorème de Huygens : $I = I_{\text{intra}} + I_{\text{inter}}$
\end{itemize}

\column{0.49\textwidth}
\begin{block}{À venir}
\textbf{Chapitre 2} : algèbre linéaire pour l'analyse des données
\begin{itemize}\small
  \item Théorème spectral
  \item Décomposition en valeurs singulières (SVD)
  \item Théorème d'Eckart-Young
  \item Projections orthogonales
\end{itemize}
\end{block}

\begin{alertblock}{Idée directrice}
Toute l'analyse des données est une
\textbf{décomposition d'inertie} via la diagonalisation
d'une matrice symétrique.
\end{alertblock}
\end{columns}
\end{frame}
